\documentclass[12pt,a4paper]{article}

% 数学相关宏包
\usepackage{amsmath, amssymb, amsthm}
\usepackage{amsfonts}
\usepackage{mathtools}
\usepackage{geometry}
\usepackage{graphicx}
\usepackage[bookmarks = false]{hyperref}
\usepackage{xcolor}
\usepackage{enumitem}
\usepackage[scheme=plain]{ctex}
\usepackage{xcolor}
\makeatletter
\renewcommand{\maketitle}{
    \begin{center}
        \vspace*{-1cm}  % 调整这个值来控制顶部间距
        {\LARGE\bfseries \@title \par}
        \vspace{0.2cm}  % 标题与作者间距
        {\large \@author }
%        \vspace{-0cm}  % 作者与日期间距
        {\normalsize \@date \par}
        \vspace{1cm}    % 日期与正文间距
    \end{center}
}
\makeatother

% 页面设置
\geometry{left=1.5cm,right=1.5cm,top=1.5cm,bottom=1.5cm}

% 超链接设置
\hypersetup{
    colorlinks=true,
    linkcolor=blue,
    urlcolor=blue,
    citecolor=red
}

% 定理环境定义
\newtheorem{theorem}{Theorem}[section]
\newtheorem{lemma}{Lemma}[theorem]
\newtheorem{corollary}{Corollary}[theorem]
\newtheorem{proposition}[theorem]{Proposition}
\newtheorem{definition}{Definition}[section]
\newtheorem{example}{Example}[section]
\newtheorem{remark}{Remark}[section]

\usepackage[noabbrev,nameinlink]{cleveref}
% 设置cleveref的标签名称
\crefname{theorem}{Theorem}{Theorems}
\Crefname{theorem}{Theorem}{Theorems}
\crefname{lemma}{Lemma}{Lemmas}
\Crefname{lemma}{Lemma}{Lemmas}
\crefname{corollary}{Corollary}{Corollaries}
\Crefname{corollary}{Corollary}{Corollaries}
\crefname{proposition}{Proposition}{Propositions}
\Crefname{proposition}{Proposition}{Propositions}
\crefname{definition}{Definition}{Definitions}
\Crefname{definition}{Definition}{Definitions}
\crefname{example}{Example}{Examples}
\Crefname{example}{Example}{Examples}
\crefname{remark}{Remark}{Remarks}
\Crefname{remark}{Remark}{Remarks}
% 数学符号定义
\newcommand{\R}{\mathbb{R}}
\newcommand{\N}{\mathbb{N}}
\newcommand{\Z}{\mathbb{Z}}
\newcommand{\Q}{\mathbb{Q}}
\newcommand{\C}{\mathbb{C}}
\newcommand{\prl}{pr_{\lambda}}
\newcommand{\prt}{pr_{\theta}}
\newcommand{\fracc}[1]{\frac{#1}{2}}
\newcommand{\indentedline}[1]{%
  \par\noindent\hspace*{#1em}
}
% 自定义 enumerate 样式：使用 (1), (1.1), (1.1.1) 等格式
% 第1级：(1), (2), ...
\setlist[enumerate,1]{label=(\arabic{enumi})}
% 第2级：(1.1), (1.2), ...
\setlist[enumerate,2]{label=(\arabic{enumi}.\arabic{enumii})}
% 第3级：(1.1.1), (1.1.2), ...
\setlist[enumerate,3]{label=(\arabic{enumi}.\arabic{enumii}.\arabic{enumiii})}
% 第4级（如有需要）
\setlist[enumerate,4]{label=(\arabic{enumi}.\arabic{enumii}.\arabic{enumiii}.\arabic{enumiv})}
% 文档信息
\title{Exalples  }
\author{Wei}
\date{\today}

\begin{document}
\maketitle
\begin{proposition} \label{prop:1}
Suppose that $\lambda \in X$, $\alpha_i \in \Phi_0^+$, $1 \leq i \leq k$, and $\beta, \gamma \in \Phi_1^+$. Let $w = s_{\alpha_k} s_{\alpha_{k-1}} \cdots s_{\alpha_1} \in \mathcal{W}$.

\begin{enumerate}
    \item Suppose that $\langle \lambda, \alpha_1^\vee \rangle < 0$. Then $(P_\lambda : M_{s_{\alpha_1}\lambda}) > 0$.
    
    \item Suppose that $\langle s_{\alpha_{i-1}} \cdots s_{\alpha_1} \lambda, \alpha_i^\vee \rangle < 0$ for all $i \in \{1,2,\ldots,k\}$. Then $(P_\lambda : M_{w\lambda}) > 0$.
    
    \item Suppose that $(\lambda, \beta) = 0$. Then $(P_\lambda : M_{\lambda + \beta}) > 0$.
    
    \item Suppose that $(\lambda, \beta) = 0$ and $\langle s_{\alpha_{i-1}} \cdots s_{\alpha_1} (\lambda + \beta), \alpha_i^\vee \rangle < 0$ for all $i \in \{1,2,\ldots,k\}$. Then $(P_\lambda : M_{w(\lambda + \beta)}) > 0$.
    
    \item Suppose that $(\lambda, \beta) = (\lambda + \beta, \gamma) = 0$ and $\mathrm{ht}(\beta) < \mathrm{ht}(\gamma)$. Then $(P_\lambda : M_{\lambda + \beta + \gamma}) > 0$.
    
    \item Suppose that $(\lambda, \beta) = (\lambda + \beta, \gamma) = 0$, $\mathrm{ht}(\beta) < \mathrm{ht}(\gamma)$, and $\langle s_{\alpha_{i-1}} \cdots s_{\alpha_1} (\lambda + \beta + \gamma), \alpha_i^\vee \rangle < 0$ for all $i \in \{1,2,\ldots,k\}$. Then $(P_\lambda : M_{w(\lambda + \beta + \gamma)}) > 0$.
\end{enumerate}
\end{proposition}


\section{gl(2|2)}
\subsection{degree of atypicality = 1}

\begin{theorem}
  Let $a,b,c \in \mathbb{Z}$ and $\lambda = (a,b|c,b)$. The projective objects $P_{\lambda}$ in $\mathfrak{B}_{a|c}$ have the following Verma flag formulae.
\end{theorem}
\begin{proof} The Cases are as follows:\par
\begin{enumerate}
  \item $a>b>c$
    \begin{enumerate}
      \item $a>b+1$ \par
      Let $\mu:=\lambda + (0,1,0,0) $. 
\begin{equation}\notag
  \begin{aligned}
      P_{\mu} = M_{a,b+1|c,b}
  \end{aligned}
\end{equation}
\begin{equation}\notag
  \begin{aligned}
    Pr_{\lambda}(P_\mu \otimes V_2) =
M_{a,b|c,b}
+M_{a,b+1|c,b+1}
  \end{aligned}
\end{equation}
All items belong to $P_\lambda$ by \cref{prop:1}.
    \end{enumerate}
\end{enumerate}
\end{proof}



\end{document}



